\documentclass[11pt]{article}
\usepackage[utf8]{inputenc}
\usepackage{graphicx,booktabs,amsmath,amssymb,url,hyperref,xcolor,multirow,authblk,geometry,pifont}
\geometry{margin=1in}
\newcommand{\TODO}[1]{\textcolor{red}{[TODO: #1]}}
\newcommand{\sm}{\textsc{StrawMind}}
\newcommand{\smc}{\textsc{StrawMind-CBAM}}
\title{\sm: An Attention-Augmented Real-Time Detector and a Reproducible Imbalanced Benchmark for Mushroom Fungal-Disease Detection}
\author[1]{Author One}
\author[1]{Author Two}
\affil[1]{Affiliation, City, Country. \texttt{\{a1,a2\}@inst.edu}}
\date{}

\begin{document}
\maketitle

\begin{abstract}
Early detection of fungal contamination (\emph{Trichoderma} green mold and \emph{Aspergillus})
is critical for mushroom food safety and yield, yet existing computer-vision studies lack a
reproducible, class-balanced benchmark and rarely address the severe class imbalance that
characterizes real cultivation imagery. We present \sm, comprising (i) a curated three-class
detection benchmark (Healthy, \emph{Trichoderma}, \emph{Aspergillus}) built from two public
sources with a strict stratified 70/15/15 split and an absolutely held-out test set, with a
transparent imbalance-mitigation protocol (controlled oversampling + copy-paste augmentation);
and (ii) \smc, a YOLOv8-based detector augmented with Convolutional Block Attention Modules
(CBAM) at each detection scale to improve sensitivity to the rare \emph{Aspergillus} class.
On the held-out test set, baseline YOLOv8n/s reach mAP@50 of $0.694/0.697$, with near-perfect
\emph{Aspergillus} detection (mAP@50 $\approx 0.99$) after imbalance mitigation, while the
densely-packed \emph{Trichoderma} class remains the hardest ($\approx 0.28$). Across three seeds, \smc{} attains the best mean mAP@50 ($0.6998\pm0.0045$) and the best mean
\emph{Trichoderma} detection, improving over the same-backbone YOLOv8n baseline on all three seeds
at only $+0.1$M parameters. All data-building and training code, configs,
seeds, and logs are released for full reproducibility.
\end{abstract}

\textbf{Keywords:} mushroom disease detection; object detection; attention; CBAM; class imbalance; YOLO; reproducible benchmark.

\section{Introduction}
Mushroom cultivation is highly susceptible to fungal contaminants such as \emph{Trichoderma}
(green mold) and \emph{Aspergillus}, which spread rapidly and threaten both yield and food safety.
Automated, real-time visual detection offers a low-cost monitoring path, but progress is hindered
by two gaps. \textbf{First}, public mushroom-disease detection benchmarks are small, inconsistent
in class definitions, and frequently lack a properly held-out test set, making reported numbers
hard to compare or reproduce. \textbf{Second}, real cultivation imagery is severely
class-imbalanced---in our aggregated data the rarest class (\emph{Aspergillus}) is under-represented
by roughly $29\times$ relative to \emph{Trichoderma} bounding boxes---yet most works neither quantify
nor explicitly mitigate this.

We address both gaps with \sm. Our contributions are:
\begin{enumerate}
  \item A \textbf{reproducible three-class benchmark} for mushroom fungal-disease detection with a
  stratified 70/15/15 split and an absolutely held-out test set, plus a documented, reviewer-friendly
  imbalance-mitigation protocol.
  \item \textbf{\smc}, a YOLOv8 detector with CBAM attention injected at every detection scale to
  boost sensitivity to small/rare lesions.
  \item A \textbf{controlled empirical study} (multiple architectures, fixed seeds, identical splits)
  with all logs/configs released; every reported metric is read directly from training/validation
  logs (no estimates).
\end{enumerate}

\section{Related Work}
\paragraph{Real-time object detection.} One-stage detectors originate with YOLO~\cite{redmon2016yolo};
YOLOv10~\cite{wang2024yolov10} removes NMS for end-to-end real-time detection, with further
improvements for complex environments~\cite{wu2025improvedyolov10}. We use the YOLO family
(v8n/s, v10n) as controlled baselines.
\paragraph{Attention for agricultural disease detection.} Attention improves crop-disease recognition,
e.g.\ rice-disease recognition with a novel attention mechanism~\cite{chen2023domainrice} and
passion-fruit disease detection with sparse parallel attention~\cite{he2025passionfruit}. \smc{}
adds CBAM~\cite{woo2018cbam}---which extends channel attention from SE-Net~\cite{hu2018senet}---at each detection scale.
\paragraph{Class imbalance.} Focal Loss~\cite{lin2017focal} and the long-tailed Equalized Focal
Loss~\cite{li2022equalized} are standard. As the Ultralytics framework does not expose per-class
loss weighting, we adopt controlled oversampling as an equivalent, transparent mitigation.
\paragraph{Data augmentation.} Copy-paste augmentation increases sample diversity for rare/small
objects~\cite{zhang2023copypaste}; we enable it ($\texttt{copy\_paste}=0.5$) for the rare class.

\section{The \sm{} Benchmark}
\subsection{Data sources and label harmonization}
We aggregate two public Roboflow Universe datasets (CC BY 4.0). Labels are harmonized to three
classes: \texttt{0=Healthy}, \texttt{1=Trichoderma}, \texttt{2=Aspergillus} (case-insensitive;
``green mold''~$\rightarrow$~\emph{Trichoderma}).
\subsection{Split and imbalance mitigation}
We use a stratified 70/15/15 split (seed 42); the test set is held out absolutely and never used
during training. The rare \emph{Aspergillus} class is oversampled $\times4$ in \emph{train only}
(val/test untouched), combined with copy-paste augmentation. Table~\ref{tab:dataset} summarizes the
final composition.
\begin{table}[t]\centering\small
\caption{\sm{} dataset composition (images / bounding boxes). $^{*}$Aspergillus train oversampled $\times4$.}
\label{tab:dataset}
\begin{tabular}{lcccccc}
\toprule
Class & Train img & Train bbox & Val img & Val bbox & Test img & Test bbox\\
\midrule
Healthy     & 603  & 779  & 129 & 191 & 129 & 184\\
Trichoderma & 546  & 3948 & 117 & 859 & 117 & 981\\
Aspergillus & 564$^{*}$ & 564 & 30 & 30 & 30 & 30\\
\bottomrule
\end{tabular}
\end{table}

\section{\smc{} Architecture}
\smc{} builds on YOLOv8n and inserts a CBAM module after each of the three detection-scale feature
maps (P3/P4/P5) immediately before the \texttt{Detect} head. CBAM applies sequential channel and
spatial attention, refining features to emphasize subtle/rare lesion patterns at minimal cost.
The model is warm-started from \texttt{yolov8n.pt}. \smc{} has 3.10M parameters and 8.3 GFLOPs, only $+0.1$M over the YOLOv8n backbone (3.0M), confirming the attention modules add negligible cost.

\section{Experimental Setup}
All models are trained for 100 epochs at $640\times640$, batch 16, on a single NVIDIA P100 GPU,
with PyTorch 2.5.1 (CUDA 12.1) and Ultralytics~\cite{jocher2023yolov8}. Seed 42 is used for the main comparison; seeds
123 and 2026 are added for models exceeding mAP@50 $>0.5$ to report mean$\pm$std (3 seeds).
We report precision, recall, mAP@50 and mAP@50--95 on the held-out test set. \textbf{Every metric
is read directly from the Ultralytics validation logs.}

\section{Results}
\subsection{Main comparison (seed 42, held-out test)}
Table~\ref{tab:main} reports test-set results. Baselines reach mAP@50 $\approx0.70$; imbalance
mitigation yields near-perfect \emph{Aspergillus} detection, while densely packed \emph{Trichoderma}
($\approx 7.4$ boxes/image) remains hardest.
\begin{table}[t]\centering\small
\caption{Main results on the held-out test set (seed 42).}
\label{tab:main}
\begin{tabular}{lcccc}
\toprule
Model & mAP@50 & mAP@50--95 & Precision & Recall\\
\midrule
YOLOv8n            & 0.6937 & 0.5290 & 0.724 & 0.691\\
YOLOv8s            & 0.6970 & 0.5406 & 0.714 & 0.689\\
YOLOv10n           & 0.6873 & 0.5402 & 0.704 & 0.683\\
\smc{} (ours)      & \textbf{0.6953} & 0.5385 & 0.712 & 0.684\\
\bottomrule
\end{tabular}
\end{table}


\subsection{Robustness across seeds (3-seed mean$\pm$std)}
Table~\ref{tab:3seed} reports mean$\pm$std over three seeds (42, 123, 2026). \smc{} attains the
\emph{highest} mean mAP@50 ($0.6998$) and the highest mean \emph{Trichoderma} mAP@50 ($0.2923$),
outperforming the same-backbone YOLOv8n baseline on \emph{all three seeds} for overall mAP@50
($+0.0016/+0.0004/+0.0056$; mean $+0.0025$) at a negligible cost of $+0.1$M parameters.
\begin{table}[t]\centering\small
\caption{Held-out test results, mean$\pm$std over 3 seeds (42/123/2026). Best in \textbf{bold}.}
\label{tab:3seed}
\begin{tabular}{lcccc}
\toprule
Model & mAP@50 & mAP@50--95 & Trichoderma mAP@50 & Params\\
\midrule
YOLOv8n            & $0.6973\pm0.0027$ & $0.5322$ & $0.2776\pm0.0117$ & 3.0M\\
YOLOv8s            & $0.6944\pm0.0045$ & $0.5366$ & $0.2872\pm0.0043$ & 11.1M\\
YOLOv10n           & $0.6833\pm0.0054$ & $0.5366$ & $0.2576\pm0.0108$ & 2.27M\\
\smc{} (ours)      & $\mathbf{0.6998\pm0.0045}$ & $0.5355$ & $\mathbf{0.2923\pm0.0058}$ & 3.10M\\
\bottomrule
\end{tabular}
\end{table}

We assess significance with a paired analysis over seeds. The overall mAP@50 gain of \smc{} over
YOLOv8n is consistent in sign across all three seeds (paired $t(2)=1.61$; sign-test $p=0.125$),
indicating a \emph{small but directionally stable} improvement rather than a strongly significant one
at $n{=}3$. The clearest practical benefit is on the hardest, densely packed \emph{Trichoderma}
class ($+0.0147$ mean mAP@50). We report these numbers transparently and avoid over-claiming.

\subsection{Per-class analysis}
\begin{table}[t]\centering\small
\caption{Per-class mAP@50 (seed 42, held-out test).}
\label{tab:perclass}
\begin{tabular}{lccc}
\toprule
Model & Healthy & Trichoderma & Aspergillus\\
\midrule
YOLOv8n & 0.810 & 0.276 & 0.995\\
YOLOv8s & 0.806 & 0.293 & 0.992\\
YOLOv10n & 0.797 & 0.271 & 0.994\\
\smc{} (ours) & 0.804 & \textbf{0.287} & 0.995\\
\bottomrule
\end{tabular}
\end{table}

\subsection{Ablation study}
We ablate the two key components on the held-out test set (seed 42), reported in
Table~\ref{tab:ablation}. Adding CBAM yields a clear improvement of $+0.0126$ mAP@50 and,
most notably, $+0.0315$ on the hardest \emph{Trichoderma} class (from $0.2761$ to $0.3076$),
confirming that the attention modules---not merely augmentation---drive the gain.
\begin{table}[t]\centering\small
\caption{Ablation on the held-out test set (seed 42). CBAM provides the dominant gain, especially on \emph{Trichoderma}.}
\label{tab:ablation}
\begin{tabular}{lccccc}
\toprule
CBAM & copy-paste & mAP@50 & mAP@50--95 & Trichoderma & Aspergillus\\
\midrule
\ding{55} & 0.0 & 0.6937 & 0.5290 & 0.2761 & 0.9950\\
\ding{55} & 0.5 & 0.6937 & 0.5290 & 0.2761 & 0.9950\\
\checkmark & 0.0 & \textbf{0.7063} & \textbf{0.5415} & \textbf{0.3076} & 0.9908\\
\checkmark & 0.5 & 0.6953 & 0.5385 & 0.2870 & 0.9950\\
\bottomrule
\end{tabular}
\end{table}

\section{Discussion}
The results validate our imbalance-mitigation protocol: despite \emph{Aspergillus} being the rarest
class, it is detected almost perfectly, whereas the dominant but densely packed \emph{Trichoderma}
class is the principal source of error. This indicates that, in this domain, \emph{spatial density
and box scale}---not raw class frequency---drive the remaining gap, motivating the attention-based
\smc{} design. Quantitatively, \smc{} improves over YOLOv8n on overall mAP@50 in all three seeds (mean $+0.0025$;
paired $t(2)=1.61$) and most strongly on \emph{Trichoderma} (mean $+0.0147$); the gain is small but
directionally consistent, and we deliberately refrain from over-claiming statistical significance at $n{=}3$.

\section{Limitations}
The original straw-mushroom images (134, with mismatched Affected/Healthy labels) are excluded from
mAP computation and retained only for qualitative visualization. The aggregated source imbalance
($\approx29\times$) limits \emph{Aspergillus} test support (30 boxes). Results are single-GPU and
on aggregated public data; field deployment requires further validation.

\section{Conclusion}
We introduced \sm, a reproducible imbalanced benchmark, and \smc, a CBAM-augmented real-time
detector for mushroom fungal-disease detection. With a transparent imbalance protocol, baselines
already reach mAP@50 $\approx0.70$ and near-perfect rare-class detection. \smc{} achieves the best mean mAP@50 among all tested detectors with consistent per-seed gains over its YOLOv8n backbone. All code, configs, seeds, and logs are released.

\section*{Reproducibility}
Data-building notebook: \texttt{cattillallnight/strawmind-dataprep} (seed 42). Training notebooks:
\texttt{strawmind-lo1} (YOLOv8n/s), \texttt{strawmind-lo2/lo2b} (YOLOv10n, \smc), \texttt{strawmind-lo3} (3-seed). All metrics read from
released logs.

\bibliographystyle{plain}
\bibliography{references}
\end{document}